% 对应单元：QM-C05-S01--QM-C05-S02
% 主要来源：Shankar 练习 5.1.1--5.2.6，印刷页 pp. 152--164，PDF pp. 164--176
% 人工核验：是

\clearpage
\section{第 5.1 节教材习题与订正}
\label{exercise-set:QM-C05-S01}

\exercisemeta
  {QM-C05-S01-EX}
  {2026-09-12}
  {Shankar，Exercises 5.1.1--5.1.4；印刷页 pp.~152--157；PDF pp.~164--169}
  {连续谱换元、自由定态、Gaussian 演化与无界算符定义域}
  {四题均已完成推导、订正与概念核验}

\exercisequestion{QM-C05-S01-E01}{教材练习 5.1.1：按能量与传播方向重写传播算符}

\textbf{对应知识点：}
\relatedtopic{QM-C05-S01-T01}{自由粒子的能谱、简并与广义本征态}。

自由粒子的传播算符在动量基中为
\begin{equation}
  U(t)=\int_{-\infty}^{\infty}
  \ket{p}\bra{p}\exp\left(-\frac{\ii p^2t}{2m\hbar}\right)\dd p.
  \label{eq:QM-C05-S01-E01-start}
\end{equation}
定义 $E=p^2/(2m)$ 和方向指标 $\alpha=\pm$，并以
\begin{equation}
  \ket{E,\alpha}\equiv
  \ket{\alpha\sqrt{2mE}}
  \label{eq:QM-C05-S01-E01-branch-definition}
\end{equation}
重新标记动量本征态。把式~\eqref{eq:QM-C05-S01-E01-start} 改写为对 $E$ 与 $\alpha$ 的求和积分，并解释测度因子。

\begin{checkedsolution}
先把动量积分分成负、正两支：
\begin{equation}
  U(t)=\int_{-\infty}^{0}(\cdots)\dd p
  +\int_0^{\infty}(\cdots)\dd p.
  \label{eq:QM-C05-S01-E01-split}
\end{equation}
由 $E=p^2/(2m)$ 得
\begin{equation}
  \frac{\dd E}{\dd p}=\frac{p}{m},
  \qquad
  \abs{\frac{\dd p}{\dd E}}
  =\frac{m}{\abs p}
  =\frac{m}{\sqrt{2mE}}.
  \label{eq:QM-C05-S01-E01-jacobian}
\end{equation}
正支 $p=+\sqrt{2mE}$ 的积分方向仍为 $0\to\infty$。负支中 $p:-\infty\to0$ 对应 $E:\infty\to0$，而 $\dd p/\dd E<0$；导数负号与反向积分限共同给出正测度。合并两支：
\begin{equation}
  \boxed{
  U(t)=\sum_{\alpha=\pm}\int_0^{\infty}
  \frac{m}{\sqrt{2mE}}
  \ket{E,\alpha}\bra{E,\alpha}
  \exp\left(-\frac{\ii Et}{\hbar}\right)\dd E}.
  \label{eq:QM-C05-S01-E01-energy-propagator}
\end{equation}
因教材的 $\ket{E,\alpha}$ 只是原 $\ket p$ 的重新标记，Jacob-ian 必须保留。若改用
\begin{equation}
  \ket{E,\alpha}_{N}
  =\left(\frac{m}{\sqrt{2mE}}\right)^{1/2}
  \ket{E,\alpha},
  \label{eq:QM-C05-S01-E01-energy-normalized-ket}
\end{equation}
则它满足能量 delta 归一化，传播算符可简写为
\begin{equation}
  U(t)=\sum_{\alpha=\pm}\int_0^\infty
  \ket{E,\alpha}_{N}{}_{N}\!\bra{E,\alpha}
  \ee^{-\ii Et/\hbar}\dd E.
  \label{eq:QM-C05-S01-E01-normalized-energy-propagator}
\end{equation}
$E=0$ 处两支合并且 Jacobian 发散，是变量换元在谱端点的坐标奇性；它不产生额外的二重简并态。
\end{checkedsolution}

\exercisequestion{QM-C05-S01-E02}{教材练习 5.1.2：在位置表象中求自由粒子能量本征函数}

\textbf{对应知识点：}
\relatedtopic{QM-C05-S01-T01}{自由粒子的能谱、简并与广义本征态}。

求解
\begin{equation}
  -\frac{\hbar^2}{2m}\psi_E''(x)=E\psi_E(x)
  \label{eq:QM-C05-S01-E02-equation}
\end{equation}
在 $E>0$ 与 $E<0$ 时的一般解，并判断这些解能否成为 $L^2(\real)$ 中的非零能量本征态。

\begin{checkedsolution}
当 $E>0$ 时，令 $p=\sqrt{2mE}$，则
\begin{equation}
  \psi_E''+\frac{p^2}{\hbar^2}\psi_E=0,
  \label{eq:QM-C05-S01-E02-positive-ode}
\end{equation}
故
\begin{equation}
  \boxed{
  \psi_E(x)=A\ee^{\ii px/\hbar}+B\ee^{-\ii px/\hbar}}.
  \label{eq:QM-C05-S01-E02-positive-solution}
\end{equation}
两项分别对应动量 $+p$ 与 $-p$，体现正能量的二重简并。它们在整条实轴上不是平方可积函数，只能作为 delta 归一化的广义本征函数；合适的连续叠加才可形成 $L^2$ 波包。

当 $E<0$ 时，令
\begin{equation}
  \kappa=\frac{\sqrt{2m\abs E}}{\hbar}>0,
  \label{eq:QM-C05-S01-E02-kappa}
\end{equation}
一般解为
\begin{equation}
  \psi_E(x)=A\ee^{\kappa x}+B\ee^{-\kappa x}.
  \label{eq:QM-C05-S01-E02-negative-solution}
\end{equation}
为使 $x\to+\infty$ 时不发散必须令 $A=0$；为使 $x\to-\infty$ 时不发散又必须令 $B=0$。因此唯一平方可积解是零函数：
\begin{equation}
  \boxed{E<0\quad\Longrightarrow\quad\psi_E=0}.
  \label{eq:QM-C05-S01-E02-no-negative-state}
\end{equation}
这与 $H=P^2/(2m)\geq0$ 完全一致。
\end{checkedsolution}

\exercisequestion{QM-C05-S01-E03}{教材练习 5.1.3：Gaussian 演化的另一种求法}

\textbf{对应知识点：}
\relatedtopic{QM-C05-S01-T03}{Gaussian 波包的平移与扩散}。

在无量纲宽度下取
\begin{equation}
  g(x)=\ee^{-x^2/2},
  \qquad
  a=\frac{\ii\hbar t}{2m}.
  \label{eq:QM-C05-S01-E03-definitions}
\end{equation}
自由演化的形式算符为 $\exp(a\partial_x^2)$。教材要求把幂级数重排为闭式 Gaussian。先通过 $a^2$ 阶的直接计算，核对
\begin{equation}
  \ee^{a\partial_x^2}g(x)
  =\frac{1}{\sqrt{1+2a}}
  \exp\left[-\frac{x^2}{2(1+2a)}\right].
  \label{eq:QM-C05-S01-E03-target}
\end{equation}

\begin{checkedsolution}
逐次求导：
\begin{align}
  g''(x)&=(x^2-1)g(x),
  \label{eq:QM-C05-S01-E03-second-derivative}\\
  g''''(x)&=(x^4-6x^2+3)g(x).
  \label{eq:QM-C05-S01-E03-fourth-derivative}
\end{align}
所以直接展开算符指数得
\begin{align}
  \ee^{a\partial_x^2}g
  &=g+ag''+\frac{a^2}{2}g''''+O(a^3)
  \notag\\
  &=g\left[
  1+(x^2-1)a
  +\left(\frac{x^4}{2}-3x^2+\frac32\right)a^2
  +O(a^3)\right].
  \label{eq:QM-C05-S01-E03-direct-series}
\end{align}

再展开目标闭式。前因子为
\begin{equation}
  (1+2a)^{-1/2}=1-a+\frac32a^2+O(a^3),
  \label{eq:QM-C05-S01-E03-prefactor-series}
\end{equation}
而
\begin{align}
  \exp\left[-\frac{x^2}{2(1+2a)}\right]
  =g(x)\left[
  1+x^2a+\left(\frac{x^4}{2}-2x^2\right)a^2
  +O(a^3)\right].
  \label{eq:QM-C05-S01-E03-exponent-series}
\end{align}
两式相乘得到
\begin{equation}
  g(x)\left[
  1+(x^2-1)a
  +\left(\frac{x^4}{2}-3x^2+\frac32\right)a^2
  +O(a^3)\right],
  \label{eq:QM-C05-S01-E03-closed-series}
\end{equation}
与式~\eqref{eq:QM-C05-S01-E03-direct-series} 逐项一致。

这个计算严谨地核验到 $a^2$ 阶，但有限阶一致本身不是全阶证明。完整证明可通过教材所示的全级数重排，或直接把初始 Gaussian 代入自由传播核积分得到式~\eqref{eq:QM-C05-S01-E03-target}。导数中出现的多项式属于 Hermite 多项式结构，全级数求和后仍为宽度变为复数的 Gaussian。
\end{checkedsolution}

\exercisequestion{QM-C05-S01-E04}{教材练习 5.1.4：无界生成元的著名反例}

\textbf{对应知识点：}
\relatedtopic{QM-C05-S01-T04}{定态方程的局部性质与算符定义域}。

考虑 $L^2(\real)$ 中的分段函数
\begin{equation}
  \psi(x)=
  \begin{cases}
    \sin(\pi x/L),&\abs x\leq L/2,\\
    0,&\abs x>L/2.
  \end{cases}
  \label{eq:QM-C05-S01-E04-state}
\end{equation}
教材指出：若把
\begin{equation}
  U(t)=\exp\left(-\frac{\ii Ht}{\hbar}\right),
  \qquad H=-\frac{\hbar^2}{2m}\frac{\dd^2}{\dd x^2}
  \label{eq:QM-C05-S01-E04-formal-evolution}
\end{equation}
机械展开为 $H$ 的幂级数，会遇到严重问题。解释问题来自哪里，以及为何精确的 $U(t)\psi$ 仍然存在。

\begin{checkedsolution}
函数在 $x=\pm L/2$ 发生跳跃：区间内的边界值分别为 $-1$ 与 $1$，区间外为零。它仍有有限范数，因而
\begin{equation}
  \psi\in L^2(\real).
  \label{eq:QM-C05-S01-E04-in-l2}
\end{equation}
但弱导数含有 Dirac delta，二阶分布导数含有 delta 的导数。令 $k=\pi/L$，可形式地写成
\begin{align}
  \psi'(x)
  &=k\cos(kx)\mathbf 1_{[-L/2,L/2]}(x)
  -\delta(x+L/2)-\delta(x-L/2),
  \label{eq:QM-C05-S01-E04-first-distribution-derivative}\\
  \psi''(x)
  &=-k^2\sin(kx)\mathbf 1_{[-L/2,L/2]}(x)
  -\delta'(x+L/2)-\delta'(x-L/2).
  \label{eq:QM-C05-S01-E04-second-distribution-derivative}
\end{align}
因此 $H\psi$ 不是 $L^2$ 函数；该态不属于自由 Hamiltonian 的定义域：
\begin{equation}
  \psi\notin\mathcal D(H).
  \label{eq:QM-C05-S01-E04-not-in-domain}
\end{equation}
于是 $H\psi$、$H^2\psi$ 等项不能作为 Hilbert 空间向量逐项使用，形式幂级数
\begin{equation}
  \sum_{n=0}^{\infty}
  \frac{1}{n!}\left(-\frac{\ii t}{\hbar}\right)^nH^n\psi
  \label{eq:QM-C05-S01-E04-invalid-series}
\end{equation}
从第一阶就失去定义。

这不表示量子时间演化不存在。自伴 $H$ 的谱定理定义
\begin{equation}
  U(t)=\int\exp\left(-\frac{\ii Et}{\hbar}\right)\dd P_H(E),
  \label{eq:QM-C05-S01-E04-spectral-definition}
\end{equation}
其中指数函数在实谱上模长恒为一。因此 $U(t)$ 是有界酉算符，
\begin{equation}
  \boxed{\mathcal D(U(t))=L^2(\real)},
  \label{eq:QM-C05-S01-E04-full-domain}
\end{equation}
精确的 $U(t)\psi$ 对本题初态完全有定义。在动量表象中，这只是在其 Fourier 变换上乘以相位 $\exp[-\ii p^2t/(2m\hbar)]$，不会改变 $L^2$ 范数。

正确区别是：$U(t)\psi$ 对每个 $L^2$ 态存在；但只有 $\psi\in\mathcal D(H)$ 时，轨道 $U(t)\psi$ 才在 $t=0$ 具有由 $-\ii H\psi/\hbar$ 给出的强导数。无界生成元的指数不能由有限维矩阵指数的逐项直觉无条件替代。
\end{checkedsolution}

\textbf{检查点：}连续谱换元必须同时处理简并分支和 Jacobian；自由平面波是广义本征态而非 $L^2$ 态；有限阶级数匹配只能作为核验，不能冒充全阶证明；在无限维 Hilbert 空间中，写下 $H\psi$ 或 $H^n\psi$ 前必须先检查定义域，而谱定理定义的 $U(t)$ 可以在整个 Hilbert 空间上存在。

\clearpage
\section{第 5.2 节教材习题与订正}
\label{exercise-set:QM-C05-S02}

\exercisemeta
  {QM-C05-S02-EX}
  {2026-09-13}
  {Shankar，Exercises 5.2.1--5.2.6；印刷页 pp.~163--164；PDF pp.~175--176}
  {无限与有限方势阱、变分下界、delta 势、绝热力和束缚态阈值}
  {六题均已完成推导、订正与概念核验}

\exercisequestion{QM-C05-S02-E01}{教材练习 5.2.1：势阱突然扩张}

\textbf{对应知识点：}
\relatedtopic{QM-C05-S02-T03}{基态能量、时间演化与表象选择}。

质量为 $m$ 的粒子最初处于长度为 $L$、中心位于原点的无限深势阱基态。势阱突然关于原点对称地扩张为原来的两倍，且瞬间波函数保持不变。求随后测量得到新势阱基态的概率。

\begin{checkedsolution}
旧势阱与新势阱的基态分别为
\begin{align}
  \psi_{\mathrm{old}}(x)
  &=\begin{cases}
    \displaystyle\sqrt{\frac2L}\cos\frac{\pi x}{L},
      &\abs x<L/2,\\[4pt]
    0,&\abs x\geq L/2,
  \end{cases}
  \label{eq:QM-C05-S02-E01-old-ground}\\
  \psi_{\mathrm{new}}(x)
  &=\begin{cases}
    \displaystyle\frac1{\sqrt L}\cos\frac{\pi x}{2L},
      &\abs x<L,\\[4pt]
    0,&\abs x\geq L.
  \end{cases}
  \label{eq:QM-C05-S02-E01-new-ground}
\end{align}
突然扩张意味着 $\psi(x,0^+)=\psi_{\mathrm{old}}(x)$。新基态振幅应在两者支集的交集上积分：
\begin{align}
  c_1^{\mathrm{new}}
  &=\int_{-L/2}^{L/2}
  \psi_{\mathrm{new}}^*(x)\psi_{\mathrm{old}}(x)\dd x
  \notag\\
  &=\frac{\sqrt2}{L}
  \int_{-L/2}^{L/2}
  \cos\frac{\pi x}{2L}\cos\frac{\pi x}{L}\dd x.
  \label{eq:QM-C05-S02-E01-overlap}
\end{align}
利用
\begin{equation}
  \cos A\cos B
  =\frac12\left[\cos(A-B)+\cos(A+B)\right],
  \label{eq:QM-C05-S02-E01-product-to-sum}
\end{equation}
有
\begin{align}
  c_1^{\mathrm{new}}
  &=\frac{\sqrt2}{2L}
  \int_{-L/2}^{L/2}
  \left[
    \cos\frac{\pi x}{2L}
    +\cos\frac{3\pi x}{2L}
  \right]\dd x
  \notag\\
  &=\frac{\sqrt2}{2L}
  \left(
    \frac{2\sqrt2L}{\pi}
    +\frac{2\sqrt2L}{3\pi}
  \right)
  =\frac{8}{3\pi}.
  \label{eq:QM-C05-S02-E01-amplitude}
\end{align}
故所求概率为
\begin{equation}
  \boxed{
  P_1=\abs{c_1^{\mathrm{new}}}^2
  =\left(\frac{8}{3\pi}\right)^2
  \approx0.7205}.
  \label{eq:QM-C05-S02-E01-probability}
\end{equation}
积分可以形式上写在新区间 $[-L,L]$，但旧波函数在 $L/2<\abs x<L$ 上为零，因此实际贡献区间仍是 $[-L/2,L/2]$。
\end{checkedsolution}

\exercisequestion{QM-C05-S02-E02}{教材练习 5.2.2：变分下界与一维吸引势}

\textbf{对应知识点：}
\relatedtopic{QM-C05-S02-T03}{基态能量、时间演化与表象选择}，以及
\relatedtopic{QM-C05-S02-T02}{有限深势阱、指数尾与束缚态量子化}。

\textbf{(a)} 设 $E_0$ 是 Hamiltonian 的最低本征值。证明任意归一化态 $\ket\psi$ 均满足
\begin{equation}
  \bra\psi H\ket\psi\geq E_0,
  \label{eq:QM-C05-S02-E02-variational-question}
\end{equation}
并在基态非简并时确定等号条件。

\textbf{(b)} 设一维势在无穷远趋于零且处处为非正的局域吸引势，$V(x)=-\abs{V(x)}$。使用归一化试探态
\begin{equation}
  \psi_\alpha(x)
  =\left(\frac\alpha\pi\right)^{1/4}
  \ee^{-\alpha x^2/2},
  \qquad\alpha>0,
  \label{eq:QM-C05-S02-E02-trial-state}
\end{equation}
说明该势至少有一个束缚态。

\begin{checkedsolution}
将归一化态在正交归一能量本征基展开：
\begin{equation}
  \ket\psi=\sum_n c_n\ket n,
  \qquad
  \sum_n\abs{c_n}^2=1.
  \label{eq:QM-C05-S02-E02-energy-expansion}
\end{equation}
于是
\begin{align}
  \bra\psi H\ket\psi
  &=\sum_{m,n}c_m^*c_nE_n\braket{m}{n}
  =\sum_n\abs{c_n}^2E_n,
  \label{eq:QM-C05-S02-E02-expectation}\\
  \bra\psi H\ket\psi-E_0
  &=\sum_n\abs{c_n}^2(E_n-E_0)\geq0.
  \label{eq:QM-C05-S02-E02-variational-proof}
\end{align}
若基态非简并，则对 $n\geq1$ 有 $E_n-E_0>0$。等号迫使所有激发态系数为零，故
\begin{equation}
  \bra\psi H\ket\psi=E_0
  \quad\Longleftrightarrow\quad
  \ket\psi=\ee^{\ii\theta}\ket0.
  \label{eq:QM-C05-S02-E02-equality-condition}
\end{equation}
若基态简并，等号条件应改为 $\ket\psi$ 属于整个基态本征子空间。

对 (b)，先由 Gaussian 积分验证
\begin{equation}
  \int_{-\infty}^{\infty}\abs{\psi_\alpha(x)}^2\dd x
  =\sqrt{\frac\alpha\pi}
  \int_{-\infty}^{\infty}\ee^{-\alpha x^2}\dd x=1.
  \label{eq:QM-C05-S02-E02-normalization}
\end{equation}
两次求导给出
\begin{equation}
  \psi_\alpha''(x)
  =(\alpha^2x^2-\alpha)\psi_\alpha(x),
  \qquad
  \expect{x^2}_\alpha=\frac1{2\alpha}.
  \label{eq:QM-C05-S02-E02-gaussian-derivatives}
\end{equation}
因此
\begin{align}
  \expect T_\alpha
  &=-\frac{\hbar^2}{2m}
  \int\psi_\alpha^*\psi_\alpha''\dd x
  =\frac{\hbar^2\alpha}{4m},
  \label{eq:QM-C05-S02-E02-kinetic}\\
  \expect V_\alpha
  &=-\sqrt{\frac\alpha\pi}
  \int_{-\infty}^{\infty}
  \ee^{-\alpha x^2}\abs{V(x)}\dd x.
  \label{eq:QM-C05-S02-E02-potential}
\end{align}
假设
\begin{equation}
  0<I_0\equiv
  \int_{-\infty}^{\infty}\abs{V(x)}\dd x<\infty.
  \label{eq:QM-C05-S02-E02-localized-assumption}
\end{equation}
控制收敛定理给出
\begin{equation}
  \int\ee^{-\alpha x^2}\abs{V(x)}\dd x
  \longrightarrow I_0
  \qquad(\alpha\to0^+).
  \label{eq:QM-C05-S02-E02-dominated-limit}
\end{equation}
故试探能量满足
\begin{align}
  E(\alpha)
  &=\frac{\hbar^2\alpha}{4m}
  -\sqrt{\frac\alpha\pi}
  \int\ee^{-\alpha x^2}\abs{V(x)}\dd x
  \notag\\
  &=\frac{\hbar^2}{4m}\alpha
  -\frac{I_0}{\sqrt\pi}\sqrt\alpha
  +o(\sqrt\alpha).
  \label{eq:QM-C05-S02-E02-trial-energy}
\end{align}
当 $\alpha\to0^+$ 时，负的 $O(\sqrt\alpha)$ 项比正的 $O(\alpha)$ 项衰减更慢，因此可取充分小的 $\alpha>0$ 使 $E(\alpha)<0$。由 (a)
\begin{equation}
  E_0\leq E(\alpha)<0.
  \label{eq:QM-C05-S02-E02-negative-ground}
\end{equation}
无穷远势能取零时，负能量位于连续散射谱阈值以下；在上述局域性条件下，它对应至少一个离散束缚态。

\textbf{严格性边界。}“任意吸引势”不能在完全没有假设时照字面使用。上述证明采用了势非零、处处非正且绝对可积等充分条件；更奇异或长程的势需要相应的谱论版本。
\end{checkedsolution}

\exercisequestion{QM-C05-S02-E03}{教材练习 5.2.3：吸引 delta 势的唯一束缚态}

\textbf{对应知识点：}
\relatedtopic{QM-C05-S02-T02}{有限深势阱、指数尾与束缚态量子化}。

考虑
\begin{equation}
  V(x)=-aV_0\delta(x),
  \qquad aV_0>0.
  \label{eq:QM-C05-S02-E03-potential}
\end{equation}
求其束缚态能量与归一化波函数，并说明为何不存在其他束缚态。

\begin{checkedsolution}
束缚态能量写成
\begin{equation}
  E=-\frac{\hbar^2\kappa^2}{2m},
  \qquad\kappa>0.
  \label{eq:QM-C05-S02-E03-negative-energy}
\end{equation}
在 $x\neq0$ 处有 $\psi''=\kappa^2\psi$。删除两端无穷远处的增长解并要求 $\psi$ 在原点连续，得到
\begin{equation}
  \psi(x)=A\ee^{-\kappa\abs x}.
  \label{eq:QM-C05-S02-E03-decaying-state}
\end{equation}
为求导数的跳跃，把定态方程在 $[-\epsilon,\epsilon]$ 上积分：
\begin{equation}
  -\frac{\hbar^2}{2m}
  \int_{-\epsilon}^{\epsilon}\psi''(x)\dd x
  -aV_0\int_{-\epsilon}^{\epsilon}
  \delta(x)\psi(x)\dd x
  =E\int_{-\epsilon}^{\epsilon}\psi(x)\dd x.
  \label{eq:QM-C05-S02-E03-integrated-equation}
\end{equation}
令 $\epsilon\to0^+$，右侧趋于零，得到
\begin{equation}
  \boxed{
  \psi'(0^+)-\psi'(0^-)
  =-\frac{2maV_0}{\hbar^2}\psi(0)}.
  \label{eq:QM-C05-S02-E03-jump-condition}
\end{equation}
对式~\eqref{eq:QM-C05-S02-E03-decaying-state}，
\begin{equation}
  \psi'(0^+)=-\kappa A,
  \qquad
  \psi'(0^-)=\kappa A.
  \label{eq:QM-C05-S02-E03-one-sided-derivatives}
\end{equation}
代入跳跃条件并排除零函数后，唯一得到
\begin{equation}
  \kappa=\frac{maV_0}{\hbar^2},
  \qquad
  \boxed{E_0=-\frac{ma^2V_0^2}{2\hbar^2}}.
  \label{eq:QM-C05-S02-E03-energy}
\end{equation}
归一化条件为
\begin{equation}
  1=2\abs A^2\int_0^\infty\ee^{-2\kappa x}\dd x
  =\frac{\abs A^2}{\kappa},
  \label{eq:QM-C05-S02-E03-normalization}
\end{equation}
故可取
\begin{equation}
  \boxed{
  \psi_0(x)=\sqrt\kappa\,\ee^{-\kappa\abs x}},
  \qquad
  \kappa=\frac{maV_0}{\hbar^2}.
  \label{eq:QM-C05-S02-E03-normalized-state}
\end{equation}
无穷远衰减与原点连续性已把任意负能态限制为式~\eqref{eq:QM-C05-S02-E03-decaying-state}，而跳跃条件又只允许一个 $\kappa$，所以没有第二个束缚能级。若尝试奇态，则连续性要求 $\psi(0)=0$，两侧单指数解只能同时为零。
\end{checkedsolution}

\exercisequestion{QM-C05-S02-E04}{教材练习 5.2.4：缓慢移动墙壁时的受力}

\textbf{对应知识点：}
\relatedtopic{QM-C05-S02-T03}{基态能量、时间演化与表象选择}。

质量为 $m$ 的粒子处于长度为 $L$ 的无限深势阱第 $n$ 个本征态。墙壁缓慢移动，使粒子保持在对应的瞬时本征态。计算 $F=-\partial E_n/\partial L$，并与能量同为 $E_n$ 的经典粒子碰撞同一面墙所产生的平均力比较。

\begin{checkedsolution}
由
\begin{equation}
  E_n=\frac{\hbar^2\pi^2n^2}{2mL^2}
  \label{eq:QM-C05-S02-E04-energy}
\end{equation}
得到量子广义力
\begin{equation}
  \boxed{
  F_{\mathrm{QM}}
  =-\pdv{E_n}{L}
  =\frac{\hbar^2\pi^2n^2}{mL^3}
  =\frac{2E_n}{L}}.
  \label{eq:QM-C05-S02-E04-quantum-force}
\end{equation}
正号表示粒子沿增大 $L$ 的方向向外推墙。

经典粒子在阱内只有动能，故
\begin{equation}
  v=\sqrt{\frac{2E_n}{m}}
  =\frac{\hbar\pi n}{mL}.
  \label{eq:QM-C05-S02-E04-speed}
\end{equation}
连续两次撞击同一面墙之间走过 $2L$，所以频率为 $v/(2L)$。一次弹性反射使粒子动量改变 $-2mv$，传给墙的动量大小为 $2mv$。因此
\begin{equation}
  F_{\mathrm{cl}}
  =\frac{v}{2L}(2mv)
  =\frac{mv^2}{L}
  =\frac{2E_n}{L}.
  \label{eq:QM-C05-S02-E04-classical-force}
\end{equation}
故
\begin{equation}
  \boxed{F_{\mathrm{cl}}=F_{\mathrm{QM}}}.
  \label{eq:QM-C05-S02-E04-force-agreement}
\end{equation}
“缓慢”保证绝热跟随，才能在尺寸变化过程中始终使用同一量子数 $n$ 的 $E_n(L)$；突然移动则应改用练习 5.2.1 的投影方法。
\end{checkedsolution}

\exercisequestion{QM-C05-S02-E05}{教材练习 5.2.5：平移后的无限深势阱}

\textbf{对应知识点：}
\relatedtopic{QM-C05-S02-T01}{无限深方势阱的边界条件与离散谱}。

把长度为 $L$ 的无限深势阱取在 $0<x<L$。从阱内一般解出发，求归一化能量本征函数与能谱，并解释平移势阱为何不改变能量。

\begin{checkedsolution}
阱内一般解为
\begin{equation}
  \psi(x)=A\cos kx+B\sin kx.
  \label{eq:QM-C05-S02-E05-general-solution}
\end{equation}
由 $\psi(0)=0$ 得 $A=0$。再由
\begin{equation}
  \psi(L)=B\sin kL=0
  \label{eq:QM-C05-S02-E05-second-boundary}
\end{equation}
并排除 $B=0$ 的零函数，得到
\begin{equation}
  k_n=\frac{n\pi}{L},
  \qquad n=1,2,\ldots.
  \label{eq:QM-C05-S02-E05-wave-numbers}
\end{equation}
由于
\begin{equation}
  \int_0^L\sin^2\frac{n\pi x}{L}\dd x=\frac L2,
  \label{eq:QM-C05-S02-E05-normalization-integral}
\end{equation}
完整本征函数和能量为
\begin{align}
  \psi_n(x)
  &=\begin{cases}
    \displaystyle\sqrt{\frac2L}\sin\frac{n\pi x}{L},&0<x<L,\\[5pt]
    0,&\text{其他位置},
  \end{cases}
  \label{eq:QM-C05-S02-E05-eigenfunction}\\
  E_n
  &=\boxed{\frac{\hbar^2\pi^2n^2}{2mL^2}}.
  \label{eq:QM-C05-S02-E05-energy}
\end{align}
势阱只从 $[-L/2,L/2]$ 平移到 $[0,L]$，宽度与形状不变。平移前后的 Hamiltonian 酉等价，因此本征函数的节点位置和坐标表达改变，谱却保持不变。
\end{checkedsolution}

\exercisequestion{QM-C05-S02-E06}{教材练习 5.2.6：有限深方势阱}

\textbf{对应知识点：}
\relatedtopic{QM-C05-S02-T02}{有限深势阱、指数尾与束缚态量子化}。

考虑
\begin{equation}
  V(x)=
  \begin{cases}
    0,&\abs x\leq a,\\
    V_0,&\abs x>a.
  \end{cases}
  \label{eq:QM-C05-S02-E06-potential}
\end{equation}
推导偶、奇束缚态的量子化条件；用无量纲图解说明能级；验证 $V_0\to\infty$ 时恢复宽度 $2a$ 的无限深势阱；最后确定第一奇束缚态出现的势阱深度阈值。

\begin{checkedsolution}
束缚态满足 $0<E<V_0$。定义
\begin{equation}
  k=\frac{\sqrt{2mE}}{\hbar},
  \qquad
  \kappa=\frac{\sqrt{2m(V_0-E)}}{\hbar},
  \label{eq:QM-C05-S02-E06-k-kappa}
\end{equation}
立刻有
\begin{equation}
  k^2+\kappa^2=\frac{2mV_0}{\hbar^2}.
  \label{eq:QM-C05-S02-E06-circle-dimensional}
\end{equation}

对偶态，在 $x\geq0$ 取
\begin{equation}
  \psi(x)=
  \begin{cases}
    A\cos kx,&0\leq x\leq a,\\
    C\ee^{-\kappa x},&x>a.
  \end{cases}
  \label{eq:QM-C05-S02-E06-even-wavefunction}
\end{equation}
在 $x=a$ 处要求 $\psi$、$\psi'$ 连续：
\begin{equation}
  A\cos ka=C\ee^{-\kappa a},
  \qquad
  -Ak\sin ka=-\kappa C\ee^{-\kappa a}.
  \label{eq:QM-C05-S02-E06-even-boundary}
\end{equation}
两式相除得到
\begin{equation}
  \boxed{k\tan ka=\kappa}.
  \label{eq:QM-C05-S02-E06-even-condition}
\end{equation}
奇态在阱内改取 $A\sin kx$，同理有
\begin{equation}
  A\sin ka=C\ee^{-\kappa a},
  \qquad
  Ak\cos ka=-\kappa C\ee^{-\kappa a},
  \label{eq:QM-C05-S02-E06-odd-boundary}
\end{equation}
故
\begin{equation}
  \boxed{k\cot ka=-\kappa}.
  \label{eq:QM-C05-S02-E06-odd-condition}
\end{equation}

令
\begin{equation}
  \alpha=ka,
  \qquad
  \beta=\kappa a,
  \qquad
  z_0^2=\frac{2mV_0a^2}{\hbar^2},
  \label{eq:QM-C05-S02-E06-dimensionless}
\end{equation}
则三组条件变为
\begin{equation}
  \alpha^2+\beta^2=z_0^2,
  \qquad
  \beta=\alpha\tan\alpha,
  \qquad
  \beta=-\alpha\cot\alpha.
  \label{eq:QM-C05-S02-E06-graph-equations}
\end{equation}
$\alpha,\beta>0$，故只看第一象限。圆与偶态或奇态曲线的每个交点同时满足能量关系和边界匹配，因而给出一个离散能级
\begin{equation}
  E=\frac{\hbar^2\alpha^2}{2ma^2}.
  \label{eq:QM-C05-S02-E06-energy-alpha}
\end{equation}

当 $V_0\to\infty$ 时，$z_0\to\infty$。对固定低能态，$\alpha$ 保持有限而 $\beta\to\infty$。于是
\begin{align}
  \text{偶态：}\quad
  \tan\alpha\to+\infty
  &\quad\Longrightarrow\quad
  \alpha\to\frac{(2r+1)\pi}{2},
  &&r=0,1,\ldots,
  \label{eq:QM-C05-S02-E06-even-infinite-limit}\\
  \text{奇态：}\quad
  \cot\alpha\to-\infty
  &\quad\Longrightarrow\quad
  \alpha\to r\pi,
  &&r=1,2,\ldots.
  \label{eq:QM-C05-S02-E06-odd-infinite-limit}
\end{align}
两族交错合并为
\begin{equation}
  \alpha_n=\frac{n\pi}{2},
  \qquad
  k_n=\frac{n\pi}{2a},
  \qquad
  \boxed{E_n=\frac{\hbar^2\pi^2n^2}{8ma^2}},
  \label{eq:QM-C05-S02-E06-recovered-spectrum}
\end{equation}
即宽度 $2a$ 的无限深势阱能谱。

下面分析束缚态数目。圆的上半支为
\begin{equation}
  \beta_{\mathrm c}(\alpha)=\sqrt{z_0^2-\alpha^2}.
  \label{eq:QM-C05-S02-E06-circle-branch}
\end{equation}
在第一偶态区间 $0<\alpha<\min(z_0,\pi/2)$ 内，$\alpha\tan\alpha$ 从零开始上升。若 $z_0<\pi/2$，当 $\alpha\to z_0^-$ 时圆降至零而偶态曲线为正；若 $z_0\geq\pi/2$，当 $\alpha\to(\pi/2)^-$ 时偶态曲线发散而圆保持有限或趋零。由于两曲线连续，介值定理保证任意 $z_0>0$ 都至少有一个偶态交点。

第一奇态正值分支位于
\begin{equation}
  \frac\pi2<\alpha<\pi.
  \label{eq:QM-C05-S02-E06-first-odd-interval}
\end{equation}
圆能真正进入该区间必须满足 $z_0>\pi/2$，等价于
\begin{equation}
  \boxed{V_0>\frac{\pi^2\hbar^2}{8ma^2}}.
  \label{eq:QM-C05-S02-E06-odd-threshold}
\end{equation}
在临界值 $z_0=\pi/2$ 处，形式交点为
\begin{equation}
  \alpha=\frac\pi2,
  \qquad
  \beta=0,
  \qquad
  E=V_0=\frac{\pi^2\hbar^2}{8ma^2}.
  \label{eq:QM-C05-S02-E06-threshold-state}
\end{equation}
但 $\beta=0$ 意味着 $\kappa=0$。阱外方程成为 $\psi''=0$，其解 $C+Dx$ 除零函数外均不属于 $L^2(\real)$。所以等号只对应连续谱边缘的阈值解；真正的奇束缚态要求严格大于式~\eqref{eq:QM-C05-S02-E06-odd-threshold} 右端。教材写成大于等于时，应把等号理解为“刚到阈值”，而不是已经出现正规化束缚态。
\end{checkedsolution}

\textbf{检查点：}无限深势阱的离散化来自两端边界条件；突然改变势阱应投影旧态，缓慢改变则可绝热跟随；变分法用任意试探态给出基态能量上界；delta 势保持波函数连续但使导数跳跃；有限阱的束缚能量由指数衰减与边界匹配共同确定，阈值解不一定是可正规化束缚态。
